\chapter{动能与势能}

\section{章节主线与学习目标}

本章研究力的空间累积效应与机械能转化规律。上一章的动量定理强调力对时间的累积：
\[
    \boldsymbol I=\int \boldsymbol F\,\mathrm dt,
\]
而本章的功与能量方法强调力对位移的累积：
\[
    W=\int \boldsymbol F\cdot\mathrm d\boldsymbol r .
\]
二者都来源于牛顿第二定律，但适用对象和物理意义不同。

\subsection{本章主线}

动能与势能一章可以按如下脉络学习：
\[
    \text{功与功率}
    \longrightarrow
    \text{动能定理}
    \longrightarrow
    \text{保守力与势能}
    \longrightarrow
    \text{功能原理}
    \longrightarrow
    \text{机械能守恒定律}.
\]
本节之后将先讨论功与功率，再由牛顿第二定律推出动能定理。保守力、势能、功能原理以及机械能守恒定律将在后续章节中继续展开。

\subsection{学习目标}

学习本章前半部分，应掌握以下内容：

\begin{enumerate}
    \item 理解功的定义
    \[
        W=\int_a^b\boldsymbol F\cdot\mathrm d\boldsymbol r,
    \]
    明确“只有力沿位移方向的分量做功”；
    \item 掌握瞬时功率
    \[
        P=\boldsymbol F\cdot\boldsymbol v,
    \]
    理解功率描述做功快慢，而不是力的大小；
    \item 掌握质点动能定理
    \[
        W_{\mathrm{net}}=\Delta E_k,
    \]
    并能用它处理已知位移、求速率的问题；
    \item 区分动量方法与能量方法：冲量改变动量，功改变动能；
    \item 能根据题目条件选择牛顿第二定律、动量定理或动能定理。
\end{enumerate}

\subsection{动量方法与能量方法的选择}

若题目强调力作用的时间、碰撞过程中的平均力或动量变化，通常优先考虑动量定理：
\[
    \boldsymbol I=\Delta\boldsymbol p.
\]
若题目强调运动路径、位移、速度大小变化或不需要求时间，通常优先考虑功与动能定理：
\[
    W_{\mathrm{net}}=\Delta E_k.
\]
简单地说，出现 $\Delta t$ 时先想到冲量，出现 $\mathrm d\boldsymbol r$ 或路径长度时先想到功。
